% !PRINTED 312
\weberpage{312}
%<<CONT
\weq{3}{(1+y)^n f\left(\frac{\alpha+\beta y}{1+y}\right) = \varphi(y) = b_0\,y^n + b_1\,y^{n-1} + \cdots + b_n,}
in which $b_0$, $b_1 \ldots b_n$ are integral homogeneous functions of degree $n$ in $\alpha$ and $\beta$, and linear functions of the coefficients of $f(x)$. Now when $x$ increases from $\alpha$ to $\beta$, $y$, likewise increasing, goes from $0$ to $\infty$. Hence $f(x)$ has exactly as many roots between $\alpha$ and $\beta$ as $\varphi(y)$ has positive roots.

From the \gesperrt{Cartesian} theorem it therefore follows:

\gesperrt{IX. The number of roots of $f(x)$ between $\alpha$ and $\beta$ is as great as, or less by an even number than, the number of changes of sign in the sequence}
\[
b_0,\ b_1,\ b_2 \ldots b_n.
\]

While, then, in the \gesperrt{Budan} theorem the changes of sign have to be counted and compared in two sequences of values, the \gesperrt{Jacobi} criterion requires only the counting of the changes of sign in one sequence of values.

In the same way the \gesperrt{Newton} rule can also be transformed, if one takes into consideration the quantities
\[
B_\nu = \nu(n-\nu)\,b_\nu^2 - (\nu+1)(n-\nu+1)\,b_{\nu-1}\,b_{\nu+1}
\]
and one thus obtains the theorem:

\gesperrt{X. The number of roots of $f(x)$ between $\alpha$ and $\beta$ is as great as, or less by an even number than, the number of changes of succession in the double sequence}
\[
\begin{aligned}
&b_0,\ b_1,\ b_2 \ldots b_n,\\
&B_0,\ B_1,\ B_2 \ldots B_n.
\end{aligned}
\]

\begin{center}
\S.\,103.

Klein's geometric comparison of the various criteria.
\end{center}

In order to compare the scope of the various criteria with one another, \gesperrt{Klein} has employed a geometric interpretation, which we shall now discuss. If one does not wish to operate with spaces of more dimensions, and thereby to lose geometric intuition, one must here
%CONT>>
